\chapter{How splatting renders, and where each method limits primitive size}
\label{ch:background}

\keybox{the idea, before the mathematics}{%
A radiance field is a model of a scene that you can photograph from a new angle.
You give it a few dozen real photographs; it fits something to them; you then ask
it for a picture from a viewpoint nobody stood at.

3D Gaussian Splatting fits a few hundred thousand small translucent blobs ---
ellipsoids with a colour and an opacity. Rendering is cheap because a blob
projects to an ellipse, and drawing a few hundred thousand ellipses in depth
order is something a GPU does very well. Fitting is gradient descent: render,
compare to the photograph, nudge every blob.

Here is the problem this thesis is about. Nothing in that procedure stops a blob
from becoming smaller than anything the photographs can resolve. If a blob ends
up a tenth of a pixel across, it still gets drawn --- at full opacity, into
whichever pixel contains its centre. At the resolution you trained on this is
invisible, because the fit compensates. Change the resolution and it is not: the
sub-pixel blobs no longer land where the fit assumed, and the image acquires
energy it should not have. That is what you are seeing in
Figure~\ref{fig:ladder}, where the same model is fine at full resolution and a
bright smear at one-eighth.

The fix, in general terms, is a \emph{band-limit}: a floor on how small a blob is
allowed to be. This chapter sets out where each published method puts that floor
and what it is made of. The rest of the thesis argues that the floor should be a matrix, with a separate limit along every direction, and measures what happens when you make it one.}

\begin{figure}[htbp]
  \centering
  \includegraphics[width=0.95\linewidth]{figures/d1-aliasing.pdf}
  \caption{The failure in one picture. The orange discs are primitives; they are
  a property of the fitted model and do not change when you change the camera.
  The grid is the pixel grid, and it does. At the rate the model was fitted at,
  the marked pixel receives three primitives and the fit has set each one's
  opacity so that the total is right. Sample half as finely and the same pixel
  receives nine --- each still contributing what was correct when it had a pixel
  to itself. The pixel gets roughly three times the light it should. Repeat over
  the whole image and the reconstruction brightens and thickens as well as blurring, which is what Figure~\ref{fig:ladder} shows happening
  to a real scene.}
  \label{fig:aliasing}
\end{figure}

\section{The emission--absorption model both baselines share}
Both baselines assume the same emission--absorption model. Along a ray
$r(t)=o+td$ with density $\sigma$ and colour $c$, the rendered radiance is
\begin{equation}
  C \;=\; \int T(t)\,\sigma(r(t))\,c(r(t),d)\,\mathrm{d}t,
  \qquad T(t)=\exp\!\Big(-\!\int_0^t\!\sigma\,\mathrm{d}u\Big).
\end{equation}
Discretised, this is the $\alpha$-blending sum both methods use. 3DGS states the
equivalence between its $\alpha$-blending and NeRF's ray marching directly; the two differ in how the field is parameterised and evaluated; what is being
rendered is the same in both. This is what allows a single argument to span both.

\section{NeRF: one frequency limit, chosen by hand, for a whole scene}
NeRF stores the field in an MLP applied to a positional encoding $\gamma(x)$ with
$L$ frequency bands. The MLP's spectral bias means the encoding band sets the
finest structure representable. The authors set $L$ by hand from the input
images. The band-limit is therefore present, explicit, global, and manually
chosen --- one scalar for the entire scene.

\section{3D Gaussian Splatting: ellipsoids, the EWA projection, and no size limit}

\begin{figure}[htbp]
  \centering
  \includegraphics[width=0.99\linewidth]{figures/d6-splatting.pdf}
  \caption{Splatting, in the three stages a reader needs before anything else in
  this thesis makes sense. A primitive is an ellipsoid sitting in the world,
  carrying a colour and an opacity. To render it, the ellipsoid is projected to an ellipse on the image plane, and the EWA transform of
  Equation~\eqref{eq:ewa} is what performs that projection. Many such ellipses are then blended into each
  pixel, nearest first. Both baselines already place a limit at stage 2, on how
  small a primitive's \emph{image} may be. This thesis is about a limit at
  stage 1, on how small the primitive \emph{itself} may be, and the two are not
  interchangeable: a stage-2 limit is recomputed for every view, whereas a
  stage-1 limit is a property of the model.}
  \label{fig:splatting}
\end{figure}
3DGS represents the field as anisotropic Gaussians with mean $p_k$ and covariance
$\Sigma_k = RSS^{\!\top}R^{\!\top}$. Projection follows EWA volume splatting: with
viewing transform $W_n$ and the Jacobian $J_n$ of the affine approximation to the
projective map,
\begin{equation}
  \Sigma' \;=\; J_n W_n \,\Sigma_k\, W_n^{\!\top} J_n^{\!\top}.
  \label{eq:ewa}
\end{equation}
This pair $(J_n, W_n)$ is the object Chapter~\ref{ch:method1} reuses. No
regularisation is applied to $\Sigma_k$; nothing bounds how small a primitive
may become.

\section{Mip-Splatting: a scalar variance floor and a screen-space dilation}
Mip-Splatting adds two filters, and it matters for Chapter~\ref{ch:results} that
they are separate objects.

\paragraph{The 2D Mip filter} is a render-side reconstruction prefilter matched
to the display sampling rate. It dilates the projected covariance by $kI$ and
compensates the opacity by $\rho = \sqrt{\lvert\Sigma'\rvert / \lvert\Sigma' +
kI\rvert}$, so that widening the splat does not brighten it. Vanilla 3DGS performs the same fixed dilation, with $k = 0.3$, and applies
\emph{no} compensation at all. The proposed method does not change this filter; but the
distinction is not cosmetic, because an arm that keeps $\rho$ is anti-aliasing
better than 3DGS in screen space whatever its 3D filter does, and
\S\ref{sec:g2} reports what that cost before it was found.

\paragraph{The 3D smoothing filter} is the acquisition-side limit at issue here:
a scalar variance
\begin{equation}
  f_k^2 \;=\; 0.2\,\Big(\frac{d_{\min,k}}{f_{\max}}\Big)^{\!2},
  \qquad \Sigma_k \leftarrow \Sigma_k + f_k^2 I,
  \label{eq:mipfilter}
\end{equation}
with $d_{\min,k}$ the smallest depth at which any camera sees primitive $k$ and
$f_{\max}$ the largest focal length over cameras, taken independently. The
opacity is compensated so that total mass is preserved under the convolution,
\begin{equation}
  \alpha_k \leftarrow \alpha_k\sqrt{|\Sigma_k| \,/\, |\Sigma_k + f_k^2 I|}.
  \label{eq:opacitycomp}
\end{equation}
Equations~\eqref{eq:mipfilter} and \eqref{eq:opacitycomp} are the baseline this
thesis generalises, and Proposition~\ref{prop:reduction} recovers them exactly.
